\section{Applying Binomial Expansion to Approximate Exponential Numbers}

Did you know we can use binomial expansion to approximate exponential
numbers? This powerful technique allows us to simplify complex expressions
and make calculations easier. We will explore how to apply binomial expansion
to approximate exponential functions.

\begin{activity}{Activity 1.6: Approximating Exponential Numbers Using Binomial Expansion $(a+b)^n$}
Ensure you have your notebook, calculator and pen ready for the activity.
\begin{enumerate}
  \item Choose a number to work with. Examples include: 1.5, 3.5, 10, etc.
  \item Express your chosen number in the form $(a+b)$. For example: if you chose
  1.5, you could write it as $(1+0.5)$.
  \item Select any integer to represent $n$. Examples include: 2, 4, $-4$.
  \item Write your decimal and $n$ in the form $(a+b)^n$. For example, using 1.5
  and $n=4$, you would write $(1+0.5)^4$.
  \item Apply the binomial theorem to expand and evaluate your expression.
  \item Choose different sets of numbers and repeat the steps above. Experiment
  with various decimal values and $n$ to see how the approximations change.
  \item Use your calculator to find the actual value of your chosen decimal raised
  to the power of $n$. For example, calculate $1.5^4$.
  \item Look at the results from your binomial expansion and the calculator.
  \item Discuss:
  \begin{enumerate}
    \item[a.] How close were your approximations to the actual values?
    \item[b.] What did you observe about the accuracy of your approximations
    with different $n$ values?
  \end{enumerate}
\end{enumerate}
\end{activity}

Great work! You have effectively used binomial expansion to approximate
exponential numbers. Now, let us solve more examples together to deepen your
understanding!

\begin{example}{Example 1.21}
Expand $(1+2y)^5$ using the binomial theorem. Hence use your results to
evaluate $(1.04)^5$, correct to 4 decimal places.

\medskip\noindent\textbf{Solution}

\begin{align*}
  (1+2y)^5 &= 1 + 5(2y) + \frac{5(4)}{2\times1}(2y)^2 + \frac{5(4)(3)}{3\times2\times1}(2y)^3 \\
  &\qquad + \frac{5(4)(3)(2)}{4\times3\times2\times1}(2y)^4 + \frac{5(4)(3)(2)(1)}{5\times4\times3\times2\times1}(2y)^5
\end{align*}
\[
  = 1 + 5(2y) + 10(4y^2) + 10(8y^3) + 5(16y^4) + 1(32y^5)
\]
\[
  (1+2y)^5 = 1 + 10y + 40y^2 + 80y^3 + 80y^4 + 32y^5
\]

$(1.04)^5$ can be written as $(1+0.04)^5$. Comparing $(1+0.04)^5$ to $(1+2y)^5$:
\[
  2y = 0.04 \implies y = 0.02
\]
Substituting $y = 0.02$ into the expansion:
\[
  1 + 10(0.02) + 40(0.02)^2 + 80(0.02)^3 + 80(0.02)^4 + 32(0.02)^5
\]
\[
  = 1 + 0.2 + 0.016 + 0.00064 + 0.0000128 + 0.0000001024 = 1.2166529024
\]
\[
  \approx 1.2167 \text{ (to 4 d.p.)}
\]
Thus $(1.04)^5 \approx 1.2167$.
\end{example}

\begin{example}{Example 1.22}
Find the first five terms of the expression $\sqrt{1+3x}$. Hence find an estimate
of the value of $\sqrt{1.12}$ to 5 decimal places.

\medskip\noindent\textbf{Solution}

\[
  \sqrt{1+3x} = (1+3x)^{1/2}
\]
\begin{align*}
  (1+3x)^{1/2} &= 1 + \frac{1}{2}(3x) + \frac{\tfrac{1}{2}\left(\tfrac{1}{2}-1\right)}{2!}(3x)^2 + \frac{\tfrac{1}{2}\left(\tfrac{1}{2}-1\right)\left(\tfrac{1}{2}-2\right)}{3!}(3x)^3 \\
  &\quad + \frac{\tfrac{1}{2}\left(\tfrac{1}{2}-1\right)\left(\tfrac{1}{2}-2\right)\left(\tfrac{1}{2}-3\right)}{4!}(3x)^4 \\[4pt]
  &= 1 + \frac{3x}{2} + \frac{-\tfrac{1}{4}}{2}(3x)^2 + \frac{\tfrac{3}{8}}{6}(3x)^3 + \frac{-\tfrac{15}{16}}{24}(3x)^4 \\[4pt]
  &= 1 + \frac{3x}{2} - \frac{1}{8}(9x^2) + \frac{1}{48}(27x^3) - \frac{5}{128}(81x^4)
\end{align*}
\[
  \sqrt{1+3x} = 1 + \frac{3x}{2} - \frac{9x^2}{8} + \frac{27x^3}{16} - \frac{405x^4}{128}
\]

$\sqrt{1.12}$ can be written as $(1+0.12)^{1/2}$. Comparing $(1+0.12)^{1/2}$ to $(1+3x)^{1/2}$:
\[
  3x = 0.12 \implies x = 0.04
\]
Substituting $x = 0.04$ into the expansion:
\[
  1 + \frac{3(0.04)}{2} - \frac{9(0.04)^2}{8} + \frac{27(0.04)^3}{16} - \frac{405(0.04)^4}{128}
\]
\[
  = 1 + \frac{0.12}{2} - \frac{0.0144}{8} + \frac{0.001728}{16} - \frac{0.00020736}{128}
\]
\[
  = 1 + 0.06 - 0.0018 + 0.000108 - 0.0000081 = 1.0582999
\]
\[
  \approx 1.05830 \text{ (to 5 d.p.)}
\]
Thus $\sqrt{1.12} \approx 1.05830$.
\end{example}

\begin{example}{Example 1.23}
Using the binomial theorem, find an estimate for the value of $\sqrt{101}$.

\medskip\noindent\textbf{Solution}

\[
  \sqrt{101} = \sqrt{100+1} = (100+1)^{1/2} = 100^{1/2}\left(1+\frac{1}{100}\right)^{1/2} = 10\left(1+0.01\right)^{1/2}
\]
\begin{align*}
  10(1+0.01)^{1/2} &= 10\left[1 + \frac{1}{2}(0.01) + \frac{\tfrac{1}{2}\left(-\tfrac{1}{2}\right)}{2!}(0.01)^2 + \frac{\tfrac{1}{2}\left(-\tfrac{1}{2}\right)\left(-\tfrac{3}{2}\right)}{3!}(0.01)^3 + \cdots\right] \\
  &= 10\left[1 + \frac{1}{2}(0.01) - \frac{1}{8}(0.0001) + \frac{1}{16}(0.000001)\right] \\
  &= 10\left[1 + 0.005 - 0.0000125 + 0.0000000625\right] \\
  &= 10\left[1.004987563\right] = 10.04987563
\end{align*}
Thus $\sqrt{101} \approx 10.0499$.
\end{example}

\begin{example}{Example 1.24}
Use the binomial theorem to find approximations for:
\begin{enumerate}
  \item[a.] $(3.97)^{1/2}$
  \item[b.] $(0.994)^{1/4}$
  \item[c.] $\sqrt{25.05}$
\end{enumerate}

\medskip\noindent\textbf{Solution}

\textbf{a.}
\[
  (3.97)^{1/2} = (4-0.03)^{1/2} = 4^{1/2}(1-0.0075)^{1/2} = 2(1-0.0075)^{1/2}
\]
\begin{align*}
  2(1-0.0075)^{1/2} &= 2\left[1 + \frac{1}{2}(-0.0075) + \frac{\tfrac{1}{2}\left(-\tfrac{1}{2}\right)}{2!}(-0.0075)^2 + \cdots\right] \\
  &= 2\left[1 - 0.00375 - \frac{1}{8}(0.0075)^2 + \cdots\right] \\
  &= 2\left[1 - 0.00375 - 0.00003125 + \cdots\right] \\
  &= 2\left[0.99621875\right] = 1.9924375
\end{align*}
\[
  \therefore (3.97)^{1/2} \approx 1.992 \text{ (to 4 s.f.)}
\]

\textbf{b.}
\[
  (0.994)^{1/4} = (1-0.006)^{1/4}
\]
\begin{align*}
  (1-0.006)^{1/4} &= 1 + \frac{1}{4}(-0.006) + \frac{\tfrac{1}{4}\left(-\tfrac{3}{4}\right)}{2!}(-0.006)^2 + \cdots \\
  &= 1 - 0.0015 - \frac{1}{32}(0.000036) + \cdots \\
  &= 1 - 0.0015 - 0.000001125 + \cdots \\
  &= 0.998498875
\end{align*}
\[
  \therefore (0.994)^{1/4} \approx 0.9985 \text{ (to 4 s.f.)}
\]

\textbf{c.}
\[
  \sqrt{25.05} = (25+0.05)^{1/2} = 25^{1/2}(1+0.002)^{1/2} = 5(1+0.002)^{1/2}
\]
\begin{align*}
  5(1+0.002)^{1/2} &= 5\left[1 + \frac{1}{2}(0.002) + \frac{\tfrac{1}{2}\left(-\tfrac{1}{2}\right)}{2!}(0.002)^2 + \cdots\right] \\
  &= 5\left[1 + 0.001 - \frac{1}{8}(0.000004) + \cdots\right] \\
  &= 5\left[1 + 0.001 - 0.0000005\right] = 5\left[1.0009995\right] \\
  &= 5.0049975
\end{align*}
\[
  \therefore \sqrt{25.05} \approx 5.005 \text{ (to 4 s.f.)}
\]
\end{example}

\subsection*{Extended Reading}
\begin{itemize}
  \item Cambridge Additional Mathematics by Michael Haese, Sandra Haese, Mark
  Humphries, Chris Sangwin. Haese Mathematics (2014).
  \item Goldie, S. (2012). \emph{Pure Mathematics}. Hodder Education.
  \item Talbert, J. F., \& Heng, H. H. (1995). \emph{Additional Mathematics: Pure \& Applied}. Pearson Education South Asia.
\end{itemize}
